\subsection{UPPER AND LOWER BOUNDS}

DEFINITION 5. Let E be a preordered set and let X be a subset of E. Any element \(x \in \mathbf{E}\) such that \(x \leqslant y\) (resp. \(x \geqslant y\) ) for all \(y \in \mathbf{X}\) is called a lower (resp. upper) bound of X in E.

Every upper bound of X is a lower bound of X with respect to the opposite ordering, and vice versa.

If \(x\) is a lower bound of \(\mathbf{X}\) , every element \(z \leqslant x\) is also a lower bound of \(\mathbf{X}\) . A lower bound of \(\mathbf{X}\) is also a lower bound of every subset of \(\mathbf{X}\) . An ordered set \(\mathbf{X}\) has a least element if and only if there exists a lower bound of \(\mathbf{X}\) which belongs to \(\mathbf{X}\) .

The set of lower bounds of a subset X of a preordered set E may be empty: this is the case when X = E and E is an ordered set which has no least element.

A subset X of E whose set of lower (resp. upper) bounds is not empty is said to be bounded below (resp. bounded above). A subset which is bounded both below and above is said to be bounded. If X is bounded below (resp. bounded above, bounded), the same is true of every subset of X.

Every subset consisting of a single element is bounded. But a subset consisting of two elements need not be bounded either above or below (no. 10).

Let E be a preordered set and let f be a mapping of an arbitrary set A into E. The mapping f is said (by abuse of language) to be bounded below (resp. bounded above, bounded) if the set \(f(\mathbf{A})\) is bounded below (resp. bounded above, bounded) in E.

